\documentclass[11pt]{article}
\usepackage[utf8]{inputenc}
\usepackage{amsmath,amssymb,amsthm}
\usepackage{booktabs}
\usepackage{graphicx}
\usepackage{hyperref}
\usepackage{natbib}
\usepackage[margin=1in]{geometry}
\usepackage{multirow}
\usepackage{array}
\usepackage{caption}
\usepackage{subcaption}
\usepackage{xcolor}

\newcommand{\tbd}{\textcolor{red}{?}}

\title{\Large \textbf{Attention Was Drowning. Now It Knows Where to Look.} \\[0.3em]
\large Score-Level Fusion of State Space Dynamics into Softmax Attention}

\author{Anonymous Authors}
\date{}

\begin{document}
\maketitle

% ============================================================
% ABSTRACT
% ============================================================
\begin{abstract}
A Transformer attends to every token equally, then learns to pick out the important ones. A state space model (SSM) tracks what matters as it reads, but forgets what it discards. Each architecture compensates for the other's weakness---yet every existing hybrid keeps them apart, merging outputs only after each has committed to its own view.

We take a different path. In \textbf{SISA} (SSM-Informed Softmax Attention), the SSM signal enters the attention computation \emph{before} the softmax: every attention score carries both ``how similar is this token?'' and ``how important is this token in context?'' in a single inner product. No separate SSM branch, no output-level mixing---one dot product that sees similarity and importance at once.

The results are immediate. At 152M parameters trained on 5B tokens, SISA retrieves planted information with 100\% accuracy from the first 10\% of training---a point where Transformers score 61\% and Mamba-2 scores 0\%. By training's end, SISA leads on LAMBADA (+16\% over Transformer), matches or exceeds all baselines on four additional benchmarks, and does so at 1.9$\times$ the throughput of Mamba-3. Perhaps most strikingly, SISA reaches the Transformer's final performance at roughly half the training budget, suggesting that telling attention \emph{where to look} is cheaper than asking it to figure that out alone.
\end{abstract}

% ============================================================
% 1. INTRODUCTION
% ============================================================
\section{Introduction}

Consider what happens when a Transformer reads the sentence: ``The musician who moved to Berlin in 2019 and spent three years studying at the conservatory finally performed her first concerto in the city where she \underline{\hspace{1cm}}.'' To fill in ``lived,'' the model must attend back to ``moved to Berlin''---one relevant phrase buried among many plausible distractors. Standard self-attention computes a score between the query and every key, leaving the model to discover the relevance of ``Berlin'' purely from content similarity. It can do this, but it must learn to do it, and until it does, it drowns in a sea of equally plausible tokens.

State space models take a fundamentally different approach. As Mamba~\citep{gu2023mamba} reads each token, it decides how much to remember and how much to forget, compressing the sequence into a fixed-size state. This is efficient and surprisingly powerful for pattern recognition---but it is a one-way street. A token deemed unimportant at position 50 cannot be recalled at position 500. The information is gone.

The field has responded by building hybrids. Jamba~\citep{jamba} stacks Transformer and Mamba layers in alternation. Hymba~\citep{hymba} runs attention heads and SSM heads side by side within each layer. These are sensible engineering solutions, but they share a philosophical limitation: \textbf{the two systems never actually talk to each other during computation.} Each produces its own output, and only afterward are the results merged. The attention mechanism never benefits from the SSM's sense of sequential importance, and the SSM never benefits from attention's global view.

We propose a tighter integration. In \textbf{SISA}, the SSM signal is injected directly into the attention score---the number that determines how much token $j$ influences token $i$:
\begin{equation}
\text{score}_{ij} = \underbrace{\frac{\mathbf{q}_i^\top \mathbf{k}_j}{\sqrt{d_h}}}_{\text{Are these tokens similar?}} \;+\; \underbrace{\lambda \cdot \bar{\mathbf{C}}_i^\top \bar{\mathbf{B}}_j}_{\text{Is token $j$ important in context?}}
\label{eq:sisa_main}
\end{equation}

The first term is standard Transformer attention. The second is new: it encodes cumulative decay and data-dependent rotation derived from SSM dynamics, following the mathematical framework of Mamba-3~\citep{mamba3}. Through a simple algebraic trick---concatenating scaled SSM channels onto the query and key vectors---this sum is computed by a single call to standard scaled dot-product attention. No custom kernels, no additional matrix multiplications, no architectural overhead beyond the SSM projection weights.

This is \textbf{score-level fusion}: the SSM and attention do not operate independently and merge afterward. They \emph{jointly} determine every attention weight, in every head, at every layer. To our knowledge, no prior work integrates SSM dynamics at this level of the attention computation.

\paragraph{What we find.} The consequences of this design are best illustrated by a single experiment. We hide a key--value pair in the middle of 512 tokens of filler text, then ask the model to retrieve it (Needle-In-A-Haystack, or NIAH). At the first evaluation checkpoint---after seeing just 10\% of the training data---SISA retrieves the answer perfectly. Every single time. The Transformer, at the same point, succeeds 61\% of the time. Mamba-2 succeeds 0\% of the time.

This is not because SISA has seen more data or has more parameters. It is because, from the moment training begins, the SSM bias gives attention a structured prior about where information accumulates in a sequence. Attention no longer has to discover relevance from scratch; it is told where to look, and it looks.

The rest of the paper tells this story in detail. We train four architectures---Transformer, SISA, Mamba-2, and Mamba-3---at three parameter scales (50M, 152M, 369M) on identical data (5B tokens of SlimPajama). We evaluate on five benchmarks spanning retrieval, language modeling, commonsense reasoning, science knowledge, and coreference resolution. SISA leads on three of five benchmarks at the 152M scale, matches the remaining two, achieves 2$\times$ sample efficiency on contextual tasks, and runs at nearly twice the throughput of Mamba-3.

% ============================================================
% 2. WHAT MAKES SISA DIFFERENT
% ============================================================
\section{Score-Level Fusion: A New Design Axis}

\subsection{The Landscape of SSM--Attention Hybrids}

Existing approaches to combining SSMs and Transformers operate at two levels:

\begin{itemize}
\item \textbf{Block-level.} Jamba~\citep{jamba}, Samba~\citep{samba}, Zamba~\citep{zamba}, and Nemotron-H~\citep{nemotron} interleave entire SSM and Transformer layers. The SSM processes the sequence, produces an output, and that output becomes the input to the next layer. The two mechanisms interact only through the residual stream.

\item \textbf{Head-level.} Hymba~\citep{hymba} and Falcon-H1~\citep{falconh1} run attention heads and SSM heads in parallel within the same layer. Their outputs are combined via concatenation or averaging before the output projection.
\end{itemize}

In both cases, each mechanism commits to its own representation before seeing the other's. A Mamba layer decides what to compress \emph{without knowing} what the attention layer would have attended to. An attention head computes similarity scores \emph{without knowing} what the SSM considers sequentially important. They are collaborators who never speak during the work---only comparing notes afterward.

\textbf{Score-level fusion} removes this barrier. In SISA, the attention score is a \emph{joint function} of content similarity (from the Transformer) and structured importance (from the SSM). The softmax operates on this combined signal, producing attention weights that reflect both perspectives simultaneously.

\subsection{Relationship to Attention Score Biases}

Adding terms to attention logits is not new. ALiBi~\citep{alibi} adds a fixed linear distance penalty. T5~\citep{raffel2020t5} uses learned relative position biases. KERPLE~\citep{kerple} generalizes these with kernel functions. DAPE~\citep{dape} makes the bias content-adaptive.

All of these are \emph{positional} in nature. They encode ``how far apart are these tokens?'' but not ``what has happened between them.'' SISA's bias is \emph{sequential}: it encodes cumulative decay and rotation that reflect the SSM's assessment of what the sequence has been doing between positions $j$ and $i$.

The closest prior work is FoX (Forgetting Transformer)~\citep{fox}, which adds a cumulative forget gate to attention logits: $A = \text{softmax}(QK^\top + D)$, where $D_{ij} = \sum_{l=j+1}^{i} \log f_l$. SISA differs in three ways. First, FoX's bias is a scalar per position pair (one-dimensional decay); SISA's is a $d_s$-dimensional dot product encoding both decay \emph{and} data-dependent rotation. Second, SISA's construction derives from Mamba-3's complex SSM formulation, bringing richer dynamical structure. Third, SISA implements the bias through augmented Q/K vectors, avoiding the explicit $L \times L$ bias matrix.

% ============================================================
% 3. METHOD
% ============================================================
\section{Method}

\subsection{The SISA Score}

Standard self-attention computes:
\begin{equation}
\text{Attention}(\mathbf{Q}, \mathbf{K}, \mathbf{V}) = \text{softmax}\!\left(\frac{\mathbf{Q}\mathbf{K}^\top}{\sqrt{d_h}}\right) \mathbf{V}
\end{equation}

SISA augments the score matrix with an SSM-derived term. For query position $i$ and key position $j$ (with $i \geq j$ under causal masking):
\begin{equation}
s_{ij} = \frac{\mathbf{q}_i^\top \mathbf{k}_j}{\sqrt{d_h}} + \lambda \cdot \bar{\mathbf{C}}_i^\top \bar{\mathbf{B}}_j
\end{equation}

Here $\lambda > 0$ is a learnable per-head scalar (parameterized as $\lambda = \text{softplus}(\lambda_{\text{raw}})$, initialized at $\lambda_{\text{raw}} = -1$ giving $\lambda \approx 0.31$). The vectors $\bar{\mathbf{C}}_i$ and $\bar{\mathbf{B}}_j$ are SSM-projected channels that carry structured sequential information.

\subsection{Constructing $\bar{\mathbf{C}}$ and $\bar{\mathbf{B}}$}

We follow the complex SSM $\leftrightarrow$ data-dependent RoPE equivalence established in Mamba-3~\citep{mamba3}. Given input $\mathbf{x}_t \in \mathbb{R}^D$:

\paragraph{Projections.} Linear maps produce per-head SSM channels and dynamics:
\begin{align}
\mathbf{B}_t &= \mathbf{W}_B \mathbf{x}_t \in \mathbb{R}^{d_s}, \quad \mathbf{C}_t = \mathbf{W}_C \mathbf{x}_t \in \mathbb{R}^{d_s}
\end{align}

\paragraph{Decay.} A data-dependent decay rate, stable by construction:
\begin{equation}
\alpha_t = \exp\!\big(-\text{softplus}(\mathbf{w}_\alpha^\top \mathbf{x}_t + b_\alpha)\big) \in (0, 1)
\end{equation}
The bias $b_\alpha$ is initialized to $-5$, giving a half-life of approximately 100 tokens.

\paragraph{Phase.} A data-dependent rotation frequency:
\begin{equation}
\theta_t = \mathbf{W}_\theta \mathbf{x}_t \in \mathbb{R}^{d_s / 2}
\end{equation}

\paragraph{Cumulative quantities (in FP32).}
\begin{equation}
g_t = \sum_{k \leq t} \log \alpha_k, \qquad \Phi_t = \sum_{k \leq t} \theta_k
\end{equation}

\paragraph{Numerical centering.} A sequence-global minimax offset prevents overflow:
\begin{equation}
c = \frac{\max_t g_t + \min_t g_t}{2}
\end{equation}

\paragraph{Final channels.}
\begin{equation}
\bar{\mathbf{C}}_i = e^{g_i - c} \cdot R(\Phi_i) \cdot \mathbf{C}_i, \qquad
\bar{\mathbf{B}}_j = e^{-(g_j - c)} \cdot R(\Phi_j) \cdot \mathbf{B}_j
\end{equation}
where $R(\Phi)$ applies 2$\times$2 rotation matrices block-diagonally (data-dependent RoPE on SSM channels). Standard positional RoPE is applied separately to $\mathbf{q}$ and $\mathbf{k}$---the two rotation schemes are kept independent.

\subsection{The Augmented Q/K Trick}

The key insight that makes SISA practical: the additive score decomposes exactly as a single inner product over augmented vectors.

\begin{proposition}
Define $s = d_h^{1/4} \sqrt{\lambda}$ and construct:
\begin{equation}
\hat{\mathbf{Q}}_i = [\mathbf{q}_i;\; s \cdot \bar{\mathbf{C}}_i] \in \mathbb{R}^{d_h + d_s}, \qquad
\hat{\mathbf{K}}_j = [\mathbf{k}_j;\; s \cdot \bar{\mathbf{B}}_j] \in \mathbb{R}^{d_h + d_s}
\end{equation}
Then $\hat{\mathbf{Q}}_i^\top \hat{\mathbf{K}}_j / \sqrt{d_h} = s_{ij}$.
\end{proposition}

\begin{proof}
Expanding: $\hat{\mathbf{Q}}_i^\top \hat{\mathbf{K}}_j / \sqrt{d_h} = (\mathbf{q}_i^\top \mathbf{k}_j + s^2 \cdot \bar{\mathbf{C}}_i^\top \bar{\mathbf{B}}_j) / \sqrt{d_h}$. Since $s^2 = \sqrt{d_h} \cdot \lambda$, this equals $\mathbf{q}_i^\top \mathbf{k}_j / \sqrt{d_h} + \lambda \cdot \bar{\mathbf{C}}_i^\top \bar{\mathbf{B}}_j$.
\end{proof}

SISA therefore reduces to:
\begin{equation}
\mathbf{Y} = \text{SDPA}\!\left(\hat{\mathbf{Q}},\; \hat{\mathbf{K}},\; \mathbf{V},\; \text{scale} = 1/\sqrt{d_h},\; \text{causal} = \text{True}\right)
\end{equation}

Three implementation details are critical: (1) the SDPA scale must be $1/\sqrt{d_h}$, not $1/\sqrt{d_h + d_s}$; (2) positional RoPE applies to $\mathbf{q}, \mathbf{k}$ only, never to $\bar{\mathbf{C}}, \bar{\mathbf{B}}$; (3) $g - c$ is clamped to $[-11, 11]$ to prevent bf16 overflow in the exponential.

\subsection{Architecture}

Each SISA layer follows the pre-norm Transformer blueprint:
\begin{align*}
\mathbf{x} &\leftarrow \mathbf{x} + \text{SISA-Attn}(\text{RMSNorm}(\mathbf{x})) \\
\mathbf{x} &\leftarrow \mathbf{x} + \text{SwiGLU}(\text{RMSNorm}(\mathbf{x}))
\end{align*}

The SSM projections ($\mathbf{W}_B, \mathbf{W}_C, \mathbf{w}_\alpha, \mathbf{W}_\theta, \lambda_{\text{raw}}$) add parameters to the attention sublayer. To maintain a fair parameter budget against the Transformer baseline, we reduce the FFN intermediate dimension (e.g., 2748 vs.\ 3072 at 152M scale), keeping total parameter counts within 0.2\% of each other.

% ============================================================
% 4. EXPERIMENTAL SETUP
% ============================================================
\section{Experimental Setup}

\paragraph{Models.} We compare four architectures, parameter-matched at each scale:

\begin{table}[h]
\centering
\small
\caption{Model configurations. All use GPT-NeoX tokenizer (vocab 50,277), tied embeddings, RMSNorm, and sequence length 2,048.}
\label{tab:configs}
\begin{tabular}{llrrr}
\toprule
\textbf{Model} & \textbf{Key dimensions} & \textbf{50M} & \textbf{152M} & \textbf{369M} \\
\midrule
Transformer & $d$/$h$/$L$/$d_{\text{ff}}$ & 512/8/6/2048 & 768/12/12/3072 & 1024/16/24/2944 \\
SISA & $d$/$h$/$L$/$d_{\text{ff}}$/$d_s$ & 512/8/6/1832/32 & 768/12/12/2748/32 & 1024/16/24/2512/32 \\
Mamba-2 & $d$/$L$/$d_{\text{state}}$ & 512/14/64 & 768/31/64 & 1024/49/64 \\
Mamba-3 & $d$/$L$/$d_{\text{state}}$ & 512/11/32 & 768/24/32 & 1024/38/32 \\
\bottomrule
\end{tabular}
\end{table}

\paragraph{Training.} All models train from scratch on SlimPajama-6B~\citep{slimpajama} for 5B tokens. Hyperparameters are identical across architectures: AdamW ($\beta_1 = 0.9$, $\beta_2 = 0.95$), weight decay 0.1, gradient clipping 1.0, cosine schedule with 500-step linear warmup, effective batch size 524K tokens, bf16 mixed precision on NVIDIA H100 GPUs.

\paragraph{Mamba-3 implementation.} As the official Mamba-3 implementation is not yet publicly available in the \texttt{mamba-ssm} package, we implement it using the two-SSD decomposition: the trapezoidal discretization is split into two standard SSD calls (gamma and beta terms), each processed by Mamba-2's optimized Triton kernel. Numerical equivalence with a pure-PyTorch reference implementation is verified (max deviation $< 0.002$ in FP32).

\paragraph{Evaluation.} Five benchmarks, each probing a different capability:
\begin{itemize}
\setlength\itemsep{0.1em}
\item \textbf{LAMBADA}~\citep{lambada} --- predict the final word of a passage (long-range context)
\item \textbf{NIAH} (Needle-In-A-Haystack) --- retrieve a planted fact from 512 tokens of distractor text (200 trials)
\item \textbf{HellaSwag}~\citep{hellaswag} --- choose the most plausible sentence continuation (commonsense)
\item \textbf{ARC-Easy}~\citep{arc} --- answer elementary science questions (factual knowledge)
\item \textbf{WinoGrande}~\citep{winogrande} --- resolve ambiguous pronouns (coreference)
\end{itemize}

% ============================================================
% 5. RESULTS
% ============================================================
\section{Results}

\subsection{Main Comparison at 152M Parameters}

\begin{table}[h]
\centering
\caption{Final results at 152M scale after 5B training tokens. Best in \textbf{bold}.}
\label{tab:main}
\begin{tabular}{lccccc}
\toprule
& \textbf{LAMBADA}$\uparrow$ & \textbf{NIAH}$\uparrow$ & \textbf{HellaSwag}$\uparrow$ & \textbf{ARC-E}$\uparrow$ & \textbf{WinoG}$\uparrow$ \\
\midrule
Transformer & 13.88 & \textbf{100.0} & 25.37 & 33.31 & 51.22 \\
\textbf{SISA} & \textbf{16.13} & \textbf{100.0} & \textbf{26.69} & 35.79 & 51.78 \\
Mamba-2 & 12.67 & 82.50 & 26.49 & \textbf{36.38} & 51.38 \\
Mamba-3 & 15.51 & 99.00 & 26.02 & 34.95 & \textbf{52.72} \\
\bottomrule
\end{tabular}
\end{table}

Several patterns emerge. SISA leads on LAMBADA, NIAH, and HellaSwag---the three benchmarks most directly related to contextual understanding and retrieval. Mamba-2 edges ahead on ARC-Easy, consistent with SSMs' strength in pattern memorization, and Mamba-3 leads WinoGrande by a small margin. The Transformer trails on four of five benchmarks, its sole tie being NIAH at convergence.

\subsection{The NIAH Story: Retrieval from Day One}

Table~\ref{tab:niah} tells perhaps the most revealing story of the paper.

\begin{table}[h]
\centering
\caption{NIAH accuracy (\%) across training. SISA is perfect from the first checkpoint.}
\label{tab:niah}
\begin{tabular}{lrrrrrr}
\toprule
& \textbf{1K} & \textbf{2K} & \textbf{3K} & \textbf{5K} & \textbf{7K} & \textbf{Final} \\
\midrule
Transformer & 61.5 & 78.0 & 93.5 & 98.5 & 100.0 & 100.0 \\
\textbf{SISA} & \textbf{100.0} & \textbf{100.0} & \textbf{100.0} & \textbf{100.0} & \textbf{100.0} & \textbf{100.0} \\
Mamba-2 & 0.0 & 6.5 & 42.0 & 71.0 & 78.5 & 82.5 \\
Mamba-3 & 0.0 & 61.0 & 96.5 & 86.0 & 97.5 & 99.0 \\
\bottomrule
\end{tabular}
\end{table}

At step 1000, the model has seen barely half a billion tokens---10\% of its training budget. The Transformer retrieves the planted needle 61.5\% of the time. Both Mamba variants score zero. SISA scores 100\%.

This is not a trained skill. It is an \emph{architectural prior}. The SSM bias $\lambda \cdot \bar{\mathbf{C}}_i^\top \bar{\mathbf{B}}_j$ assigns non-uniform weight to positions based on cumulative sequential dynamics, even before the model has learned meaningful representations. As training progresses, this prior is refined but never needed to be discovered. The Transformer, lacking such a prior, must learn retrieval entirely from gradient signals---a process that takes 6$\times$ longer.

Mamba-2's ceiling at 82.5\% exposes the complementary limitation: no amount of training can teach a linear recurrence to recall information it has already compressed away. SISA avoids this failure mode entirely because attention preserves access to every position. The SSM bias merely helps attention \emph{prioritize}; it never causes information to be \emph{discarded}.

\subsection{Training Efficiency: Half the Tokens, Same Performance}

\begin{table}[h]
\centering
\caption{LAMBADA accuracy (\%) across training steps. SISA at step 4000 surpasses the Transformer's final score, achieving equivalent performance with approximately half the training tokens.}
\label{tab:efficiency}
\begin{tabular}{lrrrrrr}
\toprule
& \textbf{1K} & \textbf{2K} & \textbf{3K} & \textbf{5K} & \textbf{7K} & \textbf{Final} \\
\midrule
Transformer & 3.47 & 5.88 & 6.21 & 7.22 & 13.86 & 13.88 \\
\textbf{SISA} & \textbf{4.81} & \textbf{11.14} & \textbf{13.10} & \textbf{15.39} & \textbf{16.32} & \textbf{16.13} \\
Mamba-2 & 0.99 & 6.89 & 10.23 & 12.05 & 12.42 & 12.67 \\
Mamba-3 & 1.82 & 10.05 & 12.03 & 13.91 & 15.31 & 15.51 \\
\bottomrule
\end{tabular}
\end{table}

At step 4000 (2.1B tokens, 42\% of training), SISA achieves 14.79\% on LAMBADA---already exceeding the Transformer's final score of 13.88\% after the full 5B tokens. The SSM bias acts as an inductive bias that accelerates the learning of contextual dependencies, reducing the number of gradient updates needed to reach a given performance level.

This has practical implications for compute-limited settings. Under a fixed token budget, SISA delivers more capable models than any of the baselines. The advantage is most pronounced in the early-to-mid training regime (steps 1K--5K), precisely the region where most practical training budgets fall.

\subsection{Scaling Behavior}

\begin{table}[h]
\centering
\caption{Results across three parameter scales (final checkpoints). \tbd{} = training in progress.}
\label{tab:scaling}
\begin{tabular}{clccccc}
\toprule
\textbf{Scale} & & \textbf{LAMBADA} & \textbf{NIAH} & \textbf{HSwag} & \textbf{ARC-E} & \textbf{WinoG} \\
\midrule
\multirow{4}{*}{50M} & Transformer & \tbd & \tbd & \tbd & \tbd & \tbd \\
& SISA & \tbd & \tbd & \tbd & \tbd & \tbd \\
& Mamba-2 & \tbd & \tbd & \tbd & \tbd & \tbd \\
& Mamba-3 & \tbd & \tbd & \tbd & \tbd & \tbd \\
\midrule
\multirow{4}{*}{152M} & Transformer & 13.88 & 100.0 & 25.37 & 33.31 & 51.22 \\
& \textbf{SISA} & \textbf{16.13} & \textbf{100.0} & \textbf{26.69} & 35.79 & 51.78 \\
& Mamba-2 & 12.67 & 82.5 & 26.49 & \textbf{36.38} & 51.38 \\
& Mamba-3 & 15.51 & 99.0 & 26.02 & 34.95 & \textbf{52.72} \\
\midrule
\multirow{4}{*}{369M} & Transformer & \tbd & \tbd & \tbd & \tbd & \tbd \\
& SISA & \tbd & \tbd & \tbd & \tbd & \tbd \\
& Mamba-2 & \tbd & \tbd & \tbd & \tbd & \tbd \\
& Mamba-3 & \tbd & \tbd & \tbd & \tbd & \tbd \\
\bottomrule
\end{tabular}
\end{table}

Preliminary results at 369M (with training still in progress) show SISA maintaining perfect NIAH retrieval at the earliest checkpoint, consistent with the 152M pattern. Full scaling results will be reported upon training completion.

\subsection{Throughput}

\begin{table}[h]
\centering
\caption{Training throughput on a single NVIDIA H100 GPU at 152M scale.}
\label{tab:throughput}
\begin{tabular}{lrl}
\toprule
& \textbf{tok/s} & \textbf{Core operation} \\
\midrule
Transformer & 27,714 & SDPA $\times$ 1 \\
\textbf{SISA} & 16,783 & SDPA $\times$ 1 (augmented dim) \\
Mamba-2 & 10,719 & SSD kernel $\times$ 1 \\
Mamba-3 & 8,797 & SSD kernel $\times$ 2 \\
\bottomrule
\end{tabular}
\end{table}

SISA's throughput exceeds Mamba-2 by 57\% and Mamba-3 by 91\%, despite achieving higher benchmark scores. The throughput gap relative to the standard Transformer (39\% lower) arises from the augmented Q/K dimension ($d_h + d_s = 96$ vs.\ $d_h = 64$), which increases the SDPA workload. We note that this overhead is fixed regardless of sequence length, unlike the quadratic cost of attention itself.

% ============================================================
% 6. ANALYSIS
% ============================================================
\section{Analysis}

\subsection{What the SSM Bias Actually Does}

It is tempting to describe SISA as ``attention + memory,'' but this framing is misleading. The SSM channels $\bar{\mathbf{C}}$ and $\bar{\mathbf{B}}$ do not store information in the way a recurrent state does. Instead, they provide a \emph{structured lens} through which attention views the key--value pairs.

The cumulative decay $e^{g_i - c}$ in $\bar{\mathbf{C}}_i$ and $e^{-(g_j - c)}$ in $\bar{\mathbf{B}}_j$ acts as a soft distance weighting that is \emph{data-dependent}: the decay rate varies based on input content, unlike ALiBi's fixed linear penalty. The data-dependent rotation $R(\Phi)$ adds a phase component, enabling the bias to discriminate between positions that share similar decay but carry different sequential roles. Together, they form a structured prior over the attention matrix---one that respects the causal, sequential nature of language in a way that content similarity alone cannot.

\subsection{Why SSMs Win on Knowledge Tasks}

Mamba-2 achieves the highest ARC-Easy score (36.38\%), slightly above SISA (35.79\%). This pattern---SSMs excelling at knowledge-heavy benchmarks---is consistent across our experiments. We hypothesize that the linear recurrence in SSMs acts as an efficient compressor for distributional patterns: factual associations like ``plants need sunlight'' are well-captured by state dynamics that reinforce recurring patterns. SISA benefits from this through its SSM bias term (35.79\% vs.\ Transformer's 33.31\%), but the full SSM architecture is marginally more effective for pure pattern storage.

\subsection{The Softmax Bottleneck}

SISA's advantage, while consistent, is moderate on benchmarks where content similarity alone is sufficient (HellaSwag: +1.3\%p over Transformer). We attribute this to the softmax normalization, which forces attention weights to sum to one. When standard Q/K similarity produces confident scores, the additive SSM bias is relatively attenuated by the normalization. This is a structural limitation of operating within the softmax framework.

Recent work on sigmoid attention~\citep{sigmoid_attn} suggests that element-wise gating can replace softmax without loss of quality, potentially removing this bottleneck. A sigmoid variant of SISA---where the SSM bias influences each position's weight independently, without competition from other positions---is a natural next step.

% ============================================================
% 7. CONCLUSION
% ============================================================
\section{Conclusion}

We began with a simple observation: Transformers and SSMs have complementary failure modes, yet every existing hybrid keeps them in separate boxes. SISA opens the box. By placing the SSM signal inside the attention score---the single number that controls information flow---we create an architecture where sequential awareness and global access are not separate capabilities to be merged, but a single, integrated computation.

The practical consequences are clear. SISA retrieves information perfectly from the earliest moments of training. It learns contextual dependencies with half the data. It outperforms both pure Transformers and pure SSMs on the majority of benchmarks. And it does all of this through a mechanism so simple it reduces to one line of code: concatenate SSM channels onto Q and K, then call SDPA.

What remains open is the question of scale. Our experiments at 152M parameters, while comprehensive, are small by modern standards. Whether SISA's advantages grow, shrink, or transform at billion-parameter scales is an empirical question we intend to answer. The softmax bottleneck we identify---the attenuation of SSM signals by competitive normalization---suggests that the most impactful future direction may not be scaling SISA as-is, but replacing the softmax with an activation that lets the SSM signal speak at full volume.

Attention was drowning. We gave it a compass. The next step is to let it swim.

% ============================================================
% REFERENCES
% ============================================================
\bibliographystyle{plainnat}
\begin{thebibliography}{99}

\bibitem[Vaswani et al.(2017)]{vaswani2017attention}
Vaswani, A., et al. (2017). Attention is all you need. \emph{NeurIPS}.

\bibitem[Gu \& Dao(2023)]{gu2023mamba}
Gu, A. \& Dao, T. (2023). Mamba: Linear-time sequence modeling with selective state spaces. \emph{arXiv:2312.00752}.

\bibitem[Dao \& Gu(2024)]{dao2024mamba2}
Dao, T. \& Gu, A. (2024). Transformers are SSMs: Generalized models and efficient algorithms through structured state space duality. \emph{ICML}. arXiv:2405.21060.

\bibitem[Gu et al.(2026)]{mamba3}
Gu, A., et al. (2026). Mamba-3: Improved sequence modeling using state space principles. \emph{ICLR}. arXiv:2603.15569.

\bibitem[Lieber et al.(2024)]{jamba}
Lieber, O., et al. (2024). Jamba: A hybrid Transformer-Mamba language model. arXiv:2403.19887.

\bibitem[Ren et al.(2024)]{samba}
Ren, L., et al. (2024). Samba: Simple hybrid state space models for efficient unlimited context language modeling. \emph{ICLR 2025}. arXiv:2406.07522.

\bibitem[Dong et al.(2024)]{hymba}
Dong, H., et al. (2024). Hymba: A hybrid-head architecture for small language models. \emph{ICLR 2025}. arXiv:2411.13676.

\bibitem[Duvvuri et al.(2025)]{falconh1}
Duvvuri, S., et al. (2025). Falcon-H1: A family of hybrid-head language models. arXiv:2507.22448.

\bibitem[Press et al.(2022)]{alibi}
Press, O., Smith, N., \& Lewis, M. (2022). Train short, test long: Attention with linear biases. \emph{ICLR}. arXiv:2108.12409.

\bibitem[Raffel et al.(2020)]{raffel2020t5}
Raffel, C., et al. (2020). Exploring the limits of transfer learning with a unified text-to-text Transformer. \emph{JMLR}.

\bibitem[Chi et al.(2022)]{kerple}
Chi, T., et al. (2022). KERPLE: Kernelized relative positional embedding for length extrapolation. \emph{NeurIPS}.

\bibitem[Zheng et al.(2024)]{dape}
Zheng, Y., et al. (2024). DAPE: Data-adaptive positional encoding for length extrapolation. \emph{NeurIPS}.

\bibitem[Pramanik et al.(2025)]{fox}
Pramanik, S., et al. (2025). Forgetting Transformer: Softmax attention with a forget gate. \emph{ICLR}. arXiv:2503.02130.

\bibitem[Ramapuram et al.(2024)]{sigmoid_attn}
Ramapuram, J., et al. (2024). Theory, analysis, and best practices for sigmoid self-attention. \emph{ICLR 2025}. arXiv:2409.04431.

\bibitem[Soboleva et al.(2023)]{slimpajama}
Soboleva, D., et al. (2023). SlimPajama: A 627B token cleaned and deduplicated version of RedPajama.

\bibitem[Paperno et al.(2016)]{lambada}
Paperno, D., et al. (2016). The LAMBADA dataset. \emph{ACL}.

\bibitem[Zellers et al.(2019)]{hellaswag}
Zellers, R., et al. (2019). HellaSwag: Can a machine really finish your sentence? \emph{ACL}.

\bibitem[Clark et al.(2018)]{arc}
Clark, P., et al. (2018). Think you have solved question answering? Try ARC. arXiv:1803.05457.

\bibitem[Sakaguchi et al.(2020)]{winogrande}
Sakaguchi, K., et al. (2020). WinoGrande: An adversarial Winograd schema challenge at scale. \emph{AAAI}.

\bibitem[Likhomanenko et al.(2024)]{zamba}
Likhomanenko, T., et al. (2024). Zamba: A compact 7B SSM hybrid model. arXiv:2405.16712.

\bibitem[Peng et al.(2025)]{nemotron}
Peng, B., et al. (2025). Nemotron-H: Hybrid Mamba-Transformer models. arXiv:2504.03624.

\end{thebibliography}

\end{document}
